\documentclass[11pt]{article}
\usepackage[a4paper,margin=1in]{geometry}
\usepackage{amsmath,amsthm,amssymb}
\usepackage{enumitem}
\usepackage{booktabs}
\usepackage{array}

\newtheorem{theorem}{Theorem}
\newtheorem{lemma}[theorem]{Lemma}
\newtheorem{corollary}[theorem]{Corollary}
\newtheorem{proposition}[theorem]{Proposition}
\theoremstyle{definition}
\newtheorem{conjecture}[theorem]{Conjecture}
\newtheorem{remark}[theorem]{Remark}
\newtheorem*{openproblem}{Open Problem / Conjecture}

\newcommand{\D}{\mathcal{D}}
\newcommand{\LB}{\mathrm{LB}}
\newcommand{\XY}{\mathrm{XY}}
\newcommand{\Sn}{S_n}

\title{IMO 2026 Problem 3 --- Solution}
\author{Proval Team}
\date{}

\begin{document}
\maketitle

\begin{center}
\textbf{Status: PARTIAL.} The value $c(n)=\dfrac{2^n}{2^{n+1}-1}$ is
\emph{conjectured} (strongly supported numerically for $n=1,\dots,5$) and is
\emph{rigorously proved} only for $n=1$ (both directions). For general $n$,
the lower bound is proved in Case~A and reduced to an open interleaving gap in
Case~B; the upper bound is open. Every gap is flagged with an
\textsc{Open Problem / Conjecture} environment. Nothing unproved is presented
as a theorem.
\end{center}

\section{Problem Statement}

Let $n$ be a positive integer. Liu Bang and Xiang Yu have a stick of length $1$
and want to divide it between themselves. Liu Bang marks at most $n$ points on
the stick, and then Xiang Yu marks at most $n$ points on the stick. The marked
points are distinct. Then, the stick is cut at all marked points, creating a
number of pieces. Afterwards, they take turns claiming any unclaimed piece of
the stick, with Liu Bang going first. Each player's goal is to maximise the
total length of their own pieces. For each $n$, determine the largest value
$c$ such that Liu Bang may guarantee a total length of at least $c$,
regardless of Xiang Yu's play.

\section{Answer}

\begin{conjecture}[Value of $c(n)$]
For every positive integer $n$,
\[
\boxed{\,c(n)=\dfrac{2^{\,n}}{2^{\,n+1}-1}\,}.
\]
Equivalently, writing $S_n:=2^{n+1}-1$, the minimax alternating-sum value
(introduced below) is $D^\star=1/S_n$, and $c(n)=\tfrac{1}{2}(1+1/S_n)=2^n/S_n$.
\end{conjecture}

\begin{table}[h]
\centering
\begin{tabular}{c|c|c|c}
\toprule
$n$ & $S_n=2^{n+1}-1$ & $D^\star=1/S_n$ & $c(n)=2^n/S_n$ \\
\midrule
$1$ & $3$  & $1/3$   & $2/3$  \\
$2$ & $7$  & $1/7$   & $4/7$  \\
$3$ & $15$ & $1/15$  & $8/15$ \\
$4$ & $31$ & $1/31$  & $16/31$\\
$5$ & $63$ & $1/63$  & $32/63$\\
\bottomrule
\end{tabular}
\caption{Verified values. $c(1)=2/3$ is proved rigorously below; $c(2),\dots,c(5)$ are confirmed by continuous optimal-XY search (numerical certificate, not a proof).}
\end{table}

\section{Setup: Reduction to the Alternating-Sum Game}

Normalise the total length to $1$. Throughout, let $S_n:=2^{n+1}-1$.

\begin{lemma}[Greedy is optimal in the alternating item-picking game]\label{lem:greedy}
Let $v_1\ge v_2\ge\cdots\ge v_m>0$ be item values. Two players pick
alternately, player~1 first, each maximising their own total. Then greedy
play---always take the largest remaining item---is a subgame-perfect
equilibrium for \emph{both} players; the resulting payoff is that player~1
receives $v_1+v_3+v_5+\cdots$ (odd positions) and player~2 receives
$v_2+v_4+v_6+\cdots$ (even positions).
\end{lemma}

\begin{proof}
By strong induction on $m$. The cases $m\le 2$ are immediate. For $m\ge 3$,
suppose player~1's first pick is $v_j$ ($j\in\{1,\dots,m\}$); write $R_j$ for
the multiset $\{v_1,\dots,v_m\}\setminus\{v_j\}$. After player~1 takes $v_j$,
player~2 moves first on $R_j$; by the induction hypothesis (applied to the
$(m-1)$-item game on $R_j$, with player~2 now in the ``first player'' role),
player~2 obtains the odd-position sum of $R_j$ (sorted) and player~1 obtains
the even-position sum of $R_j$. Hence player~1's total payoff is
\[
\mathrm{payoff}(j)=v_j+(\text{even-position sum of the sorted multiset }R_j).
\]

We claim $\mathrm{payoff}(1)\ge\mathrm{payoff}(j)$ for every $j$.

Sort $R_j$: it is $v_1,\dots,v_{j-1},v_{j+1},\dots,v_m$ (the original order
with $v_j$ removed, still decreasing). Its even-position ($2$nd, $4$th,
$\dots$) elements are:
\begin{itemize}[leftmargin=2em]
\item if $j=2k+1$ is odd: $\{v_2,v_4,\dots,v_{2k}\}\cup\{v_{2k+3},v_{2k+5},\dots\}$;
\item if $j=2k$ is even: $\{v_2,v_4,\dots,v_{2k-2}\}\cup\{v_{2k+1},v_{2k+3},\dots\}$.
\end{itemize}

In the odd case $j=2k+1$,
\begin{align*}
\mathrm{payoff}(j) &= v_{2k+1}+(v_2+v_4+\cdots+v_{2k})+(v_{2k+3}+v_{2k+5}+\cdots),\\
\mathrm{payoff}(1) &= v_1+(v_3+v_5+\cdots)\\
&= v_1+(v_3+v_5+\cdots+v_{2k-1})+v_{2k+1}+(v_{2k+3}+v_{2k+5}+\cdots).
\end{align*}
Subtracting,
\[
\mathrm{payoff}(1)-\mathrm{payoff}(j)=(v_1-v_2)+(v_3-v_4)+\cdots+(v_{2k-1}-v_{2k})\ge 0,
\]
since the sequence is decreasing. The even case $j=2k$ is analogous:
\[
\mathrm{payoff}(1)-\mathrm{payoff}(j)=(v_1-v_2)+(v_3-v_4)+\cdots+(v_{2k-1}-v_{2k})\ge 0.
\]
Thus player~1's optimal first move is $v_1$ (greedy), and the induction
hypothesis then forces greedy thereafter. The same argument, applied to
player~2 as the first mover on $R_1$, shows greedy is optimal for player~2 as
well.
\end{proof}

\begin{corollary}[Reduction]\label{cor:reduction}
In the original game, after all cuts the final pieces sorted descending
$b_1\ge b_2\ge\cdots\ge b_m$ satisfy: under optimal play in the claiming phase,
Liu Bang receives $\tfrac{1}{2}(1+\D)$ and Xiang Yu receives
$\tfrac{1}{2}(1-\D)$, where
\[
\D:=b_1-b_2+b_3-b_4+\cdots=\sum_{i}(-1)^{i+1}b_i
\]
is the alternating sum of the resulting multiset. Hence LB's guaranteed payoff
equals $\tfrac{1}{2}(1+\min_{\XY}\D)$, and the game reduces to:
\begin{itemize}[leftmargin=2em]
\item LB chooses a partition of $[0,1]$ into at most $n+1$ pieces;
\item XY then adds at most $n$ cuts;
\item the payoff-relevant quantity is the alternating sum $\D$ of the resulting
multiset; LB maximises the guaranteed $\D$, XY minimises it.
\end{itemize}
The conjectured value is $\min_{\XY}\D\ge 1/S_n$ for LB's geometric partition,
and $\min_{\XY}\D\le 1/S_n$ against every LB partition.
\end{corollary}

\begin{proof}[Proof of Corollary~\ref{cor:reduction}]
The total length is $1=\sum_i b_i$; LB's odd-position sum plus XY's
even-position sum is $1$, and their difference is $\D$, so LB's take is
$(1+\D)/2$. Substituting the guaranteed $\D=\min_{\XY}\D$ gives LB's guaranteed
payoff.
\end{proof}

Two reformulations are used in the analysis:
\begin{itemize}[leftmargin=2em]
\item[(i)] $\D=1-2\cdot(\text{XY's even-position sum})$, so minimising $\D$
is the same as maximising XY's take.
\item[(ii)] \emph{Parity-integral identity:} $\D=\int_0^{b_1}\mathbf{1}_{r(t)\text{ odd}}\,dt$,
where $r(t):=\#\{\text{pieces of size}\ge t\}$; each piece $b_i$ contributes
the interval $(0,b_i]$ to $r(t)$, and summing the parity indicator over $t$
gives the alternating sum.
\end{itemize}

\section{Lower Bound: $\D\ge 1/S_n$ for LB's Geometric Partition (Partial)}

We scale to integer units: LB plays the \emph{geometric partition}
\[
G_n=(1,2,4,\dots,2^n),\qquad\text{total }S_n=2^{n+1}-1,
\]
and the target becomes $\D\ge 1$ in these units (i.e.\ $\D\ge 1/S_n$ after
rescaling to total $1$). The central structural fact is the
\textbf{``$+1$ gap''}:
\begin{equation}\label{eq:plusone}
\text{for every }k\ge 1,\qquad 2^k=(1+2+\cdots+2^{k-1})+1. \tag{$\dagger$}
\end{equation}
Equivalently, each geometric piece exceeds the sum of ALL smaller geometric
pieces by exactly $1$; in particular the largest piece $2^n$ exceeds the sum of
all the others ($2^n-1$) by exactly $1$---the target.

Define the proposition
\[
\mathbf{L(n)}:\;\text{every refinement of }G_n\text{ by at most }n\text{ cuts has }\D\ge 1.
\]

\subsection{Base case $L(1)$ (proved)}

$G_1=(1,2)$, total $3$, one cut allowed. XY either cuts the piece $2$ or cuts
the piece $1$.
\begin{itemize}[leftmargin=2em]
\item \textbf{Cut $2$ into $(a,2-a)$, $0<a<2$.} The three pieces are
$\{1,a,2-a\}$. By symmetry ($a\leftrightarrow 2-a$) assume $a\ge 1$, so
$2-a\le 1\le a$; the sorted order is $a,1,2-a$, and
$\D=a-1+(2-a)=1$. (At $a=1$ this is the halving $2\to 1+1$, giving
$\{1,1,1\}$ with $\D=1-1+1=1$.) So cutting $2$ always gives $\D=1$.
\item \textbf{Cut $1$ into $(a,1-a)$, $0<a<1$.} The pieces are
$\{2,a,1-a\}$; sorted (assume $a\ge 1-a$, i.e.\ $a\ge 1/2$, by symmetry) as
$2,a,1-a$, and $\D=2-a+(1-a)=3-2a$. Since $0<a<1$, we have
$1<3-2a<3$; in particular $\D>1$ strictly.
\end{itemize}
In every case $\D\ge 1$, with equality iff XY halves the piece $2$. This proves
$L(1)$: the bound $1/S_1=1/3$ is attained (after rescaling) by LB's geometric
partition $(1/3,2/3)$ and XY's halving of the $2/3$-piece. $\qed$

\subsection{Case A of the inductive step (proved for all $n$)}

\begin{lemma}[Case A: XY does not cut $2^n$]\label{lem:caseA}
Suppose XY uses $0$ cuts on the piece $2^n$ (so $2^n$ is intact and all $\le n$
cuts fall on the tail $(1,2,\dots,2^{n-1})$). Then $\D\ge 1$.
\end{lemma}
\begin{proof}
Then $b_1=2^n$, and every other piece is at most $2^{n-1}<2^n$, so $2^n$
occupies position~$1$ alone. Writing $\D_{\text{tail}}$ for the alternating
sum of the refined tail (positions $2,3,\dots$, with the sign $-$ at
position~$2$),
\[
\D=2^n-\D_{\text{tail}}.
\]
Now $\D_{\text{tail}}=(\text{odd-position tail pieces})-(\text{even-position tail
pieces})\le(\text{odd-position tail pieces})\le(\text{sum of all tail pieces})
=2^n-1$. Therefore
\[
\D\ge 2^n-(2^n-1)=1. \qedhere
\]
\end{proof}

Case~A is far stronger than required: it does not even need the inductive
hypothesis, only that the total tail mass is $2^n-1$. It shows XY is forced to
cut $2^n$ if XY hopes to approach the bound.

\subsection{Case B of the inductive step (framework; GAP)}

Suppose XY uses $k\ge 1$ cuts on $2^n$, producing fragments
$F=(f_1\ge f_2\ge\cdots\ge f_{k+1}>0)$ with $\sum f_i=2^n$ (consuming $k$
cuts), and refines the tail $G_{n-1}=(1,\dots,2^{n-1})$ with at most
$n-k\le n-1$ cuts into a multiset $T$ of total $2^n-1$. The final multiset is
the sorted merge of $F$ and $T$.

\textbf{Intended inductive reduction.} The tail $T$ is a refinement of
$G_{n-1}$ by $\le n-1$ cuts, so the inductive hypothesis $L(n-1)$ applies to
$T$ \emph{as a standalone game} and gives $\D(T)\ge 1$. The intended mechanism
is the ``$+1$ gap'' invariant~\eqref{eq:plusone}: cutting the dominant piece
$2^n$ sends its fragments to positions where they contribute favourably to $\D$
\emph{unless} XY interleaves a fragment between the top ranks of the merge;
any such interleaving consumes a cut, and the remaining $\le n-k$ cuts then
face (in the relevant sub-structure) the smaller geometric instance $G_{n-1}$,
to which $L(n-1)$ applies and forces the residual alternating sum to be $\ge 1$.
The equality case is \emph{full halving}---split every geometric piece $2^j$
($j=1,\dots,n$) into two equal halves $2^{j-1}+2^{j-1}$, producing the $2n+1$
pieces $2^{n-1}(\times2),2^{n-2}(\times2),\dots,1(\times3)$, whose alternating
sum is exactly $1$ (every consecutive pair cancels, leaving a single unit).

\textbf{Settled instances.} $L(1)$ is proved above. $L(2)$ (i.e.\ $\D\ge 1$ for
$G_2=(1,2,4)$ under $\le 2$ cuts) is verified exhaustively over the continuous
cut positions: the function $\D$ is piecewise-linear in the two cut positions,
its breakpoints occur exactly where two pieces coincide in the sorted order,
and at every such breakpoint one checks $\D\ge 1$ directly (the minimum,
$\D=1$, is attained uniquely at full halving $(4\to2+2,\;2\to1+1)$, giving
$\{2,2,1,1,1\}$ with $\D=2-2+1-1+1=1$). The same continuous breakpoint check
confirms $L(3),L(4),L(5)$.

\begin{openproblem}[General lower-bound Case B]
The general Case~B requires proving, for arbitrary $n$ and arbitrary fragment
multisets $F$ of $2^n$ (arising from $k\ge 1$ cuts) merged with arbitrary
refined tails $T$ (arising from $\le n-k$ cuts on $G_{n-1}$), that the merged
alternating sum is $\ge 1$.
\begin{itemize}[leftmargin=2em]
\item The bold inequality ``$\D(\mathrm{merge})\ge\sum F-\sum T=1$'' is
\textbf{false} for arbitrary multisets (counterexample:
$F=\{5,5,5,5\}$, $T=\{4,4,4,4\}$: the sorted merge
$5,5,5,5,4,4,4,4$ has $\D=0<4=\sum F-\sum T$).
\item The merge lemma is TRUE (tested, $0$ violations in $20\,000$ trials,
$n=2,\dots,5$) when $T$ is an \emph{actual} refinement of $G_{n-1}$ (the real
inductive hypothesis), but FALSE for general tails with $\D(T)\ge 1$ (thousands
of violations, min $\D\approx 0.03$). Hence the induction must carry the
\emph{dyadic/geometric structure} of the tail (each piece is a fragment of some
$2^j$, $j\le n-1$), not merely the numeric bound $\D(T)\ge 1$.
\item A clean invariant formalising ``the $+1$ gap survives refinement under
$\le n$ cuts''---for instance a parity/rank-function argument via the identity
$\D=\int\mathbf{1}_{r(t)\text{ odd}}\,dt$, where each cut of a piece of size
$s$ into $(m,M)$ ($m\le M$, $m+M=s$) changes $r(t)$ by $+1$ on $(0,m]$ and $-1$
on $(M,s]$---has \textbf{not} been closed.
\end{itemize}
This is the \textbf{open lower-bound gap}: the interleaving of the fragments
of $2^n$ with the refined tail $G_{n-1}$, using the cut budget $\le n$ together
with~\eqref{eq:plusone}, is not proved here.
\end{openproblem}

\section{Upper Bound: $\D\le 1/S_n$ Against Every LB Partition (Partial)}

\subsection{Base case $U(1)$ (proved, boundary fixed)}

LB makes $\le 2$ pieces; let the largest be $L\ge 1/2$ and the rest have total
$R=1-L\le 1/2$. (If LB makes a single piece, $L=1,R=0$, covered by the first
branch.) XY has $1$ cut. Target $\D\le 1/3$.
\begin{itemize}[leftmargin=2em]
\item \textbf{If $L\ge 2/3$:} XY halves $L$ into $(L/2,L/2)$. The pieces are
$L/2,L/2,R$; since $L\ge 2/3$ we have $L/2\ge 1/3\ge R$, so the sorted order is
$L/2,L/2,R$, and $\D=L/2-L/2+R=R=1-L\le 1/3$. Equality at $L=2/3$ (geometric).
\item \textbf{If $1/2\le L<2/3$} (so $R=1-L>1/3$ and $2L-1<1/3$): XY shaves a
sliver $t$ off $L$, with $0<t\le\min(R,2L-1)$ (such $t$ exists since $R>1/3>0$
and $2L-1\ge 0$ when $L>1/2$), producing $(L-t,t)$. Because $t\le R$ and
$t\le 2L-1\le L-t$, the sorted order is $(L-t),R,t$, and
$\D=(L-t)-R+t=L-R=2L-1<1/3$.
\item \textbf{Boundary $L=1/2$ exactly} (so $R=1/2$, two equal pieces---the
shaving argument above degenerates since $2L-1=0$ forces $t\le 0$): XY cuts
one of the two $1/2$-pieces into $(a,1/2-a)$ with $0<a\le 1/4$ (so
$a\le 1/2-a$). The three pieces are $1/2$ (uncut), $1/2-a,a$; sorted descending
$1/2,1/2-a,a$, and $\D=1/2-(1/2-a)+a=2a$. XY chooses $a\le 1/6$, giving
$\D=2a\le 1/3$ (and $a\to 0$ drives $\D\to 0$). Hence $\D\le 1/3=1/S_1$.
\end{itemize}
All three branches give $\D\le 1/3=1/S_1$, with equality at $L=2/3$ (LB's
geometric partition) and at the boundary $L=1/2$ with the $(1/3,1/6)$ cut.
Combined with the lower bound $L(1)$, this proves $c(1)=2/3$ rigorously. $\qed$

\subsection{General $U(n)$ (open)}

\begin{openproblem}[General upper bound]
For $n\ge 2$, no rigorous XY strategy is proved here. The following routes are
recorded as dead ends or open:
\begin{itemize}[leftmargin=2em]
\item \textbf{Myopic greedy} (``always cut to minimise $\D$ immediately'')
\textbf{fails}: $36/400$ failures at $n=2$, $67/400$ at $n=3$; e.g.\ on
$(0.4,0.35,0.25)$ greedy-halve gives $\D=0.25>1/7$, whereas the optimal
``match'' gives $\D=0.05$.
\item \textbf{``Halve-or-match'' hybrid} (halve the largest $a_1$ if
$a_1\ge 2a_2$, else match $a_1$ to $a_2$) reduces the config by one piece per
cut and preserves $\D$ = alternating sum of the residual conceptual config,
giving exactly $1/S_n$ on the geometric config. But the scale-invariant
induction $V(C,m)\le T/S_m$ (total $T$, $m$ cuts) \textbf{fails for $m\ge 2$}:
the per-step reduction is $\max(a_1,2a_2)\ge 2T/(m+2)$, and
$2T/(m+2)<T\cdot 2^m/S_m$ for $m\ge 2$, so neither branch always meets its
sufficient condition. In the ``flat'' regime (both $a_1<T\cdot 2^m/S_m$ and
$a_2<T\cdot 2^{m-1}/S_m$) the hybrid under-delivers.
\item \textbf{Hall/pairing} (pair pieces into cancelling equal-consecutive
pairs, each contributing $0$ to $\D$, leaving a leftover $\le 1/S_n$) is not
forced for arbitrary LB configs: the pairing strategies tested leave a leftover
exceeding $1/S_n$ on some random configs for $n\ge 3$
($n=3$: $0.0757>1/15$; $n=4$: $0.0407>1/31$; $n=5$: $0.0228>1/63$).
\item \textbf{Majorisation / smoothing (most promising, unproved).} The
function $f(C):=\XY\text{'s optimal }\D\text{ against the LB partition }C$
appears to be Schur-maximised \emph{uniquely} at the geometric partition: for
$n=2,3$, over $200$--$800$ random configs the ratio
$f(C)/(1/S_n)$ never exceeds $1$ and equals $1$ only at the geometric config.
A smoothing lemma---every LB partition can be continuously deformed toward the
geometric partition while only \emph{increasing} $f$---would immediately yield
the upper bound from the lower-bound value $f(\text{geometric})=1/S_n$. This
smoothing monotonicity is the conjectured but \textbf{unproved} heart of the
upper bound.
\end{itemize}
\end{openproblem}

\section{Numerical Evidence}

A continuous optimal-XY search was run for $n=1,\dots,5$. The candidate cut set
is $\{\text{fragment}=\text{an existing piece}\}\cup\{\text{fragment}=\text{half a
piece}\}$, which contains every breakpoint at which the sorted order (and hence
the piecewise-linear form $\D$) changes slope; the search therefore returns the
exact piecewise-linear minimum of $\D$ as a function of the cut positions. This
is a \emph{numerical certificate}, not a proof step.

\begin{center}
\begin{tabular}{c|c|c|c}
\toprule
$n$ & $\min_{\XY}\D$ (search) & $1/S_n$ & LB partition attaining it \\
\midrule
$1$ & $1/3$   & $1/3$   & $(1/3,2/3)$ \\
$2$ & $1/7$   & $1/7$   & $(1/7,2/7,4/7)$ \\
$3$ & $1/15$  & $1/15$  & $(1/15,2/15,4/15,8/15)$ \\
$4$ & $1/31$  & $1/31$  & geometric \\
$5$ & $1/63$  & $1/63$  & geometric \\
\bottomrule
\end{tabular}
\end{center}
In every case the minimum is attained \emph{uniquely} by LB playing the
geometric partition and XY playing full halving ($2^j\to2^{j-1}+2^{j-1}$ for
each $j=1,\dots,n$). Non-geometric LB configs were searched on a grid
($n=1,2,3$) with full-lookahead optimal XY: the max guaranteed $\D$ is attained
\emph{uniquely} at the geometric config; every other config lets XY drive $\D$
strictly below $1/S_n$ (often to $0$). The conjecture is not contradicted.

\section{Status}

\begin{itemize}[leftmargin=2em]
\item \textbf{Proved rigorously:} the greedy-pick reduction
(Lemma~\ref{lem:greedy}, Corollary~\ref{cor:reduction}); the lower bound for
$n=1$ ($L(1)$); the lower-bound Case~A for all $n$ (Lemma~\ref{lem:caseA}); the
upper bound for $n=1$ ($U(1)$, all three branches incl.\ the $L=1/2$ boundary).
These together give a \textbf{complete rigorous proof of $c(1)=2/3$}.
\item \textbf{Conjectured (not proved), strongly supported numerically:}
$c(n)=2^n/(2^{n+1}-1)$ for all $n$, confirmed for $n=2,3,4,5$ by continuous
optimal-XY search.
\item \textbf{Open gap (lower bound):} general Case~B---the interleaving of the
fragments of $2^n$ with the refined tail $G_{n-1}$ under the cut budget $\le n$
and the ``$+1$ gap''~\eqref{eq:plusone}. Requires a structural invariant carrying
the dyadic structure of the tail, not merely $\D(T)\ge 1$.
\item \textbf{Open gap (upper bound):} a rigorous XY strategy for $n\ge 2$, or a
majorisation/smoothing proof that the geometric config Schur-maximises XY's
optimal $\D$.
\end{itemize}

The general case is a \textbf{conjecture} supported by numerical evidence, not a
complete proof. The status of this solution is \textbf{PARTIAL}.

\end{document}
